\textbf{Partie A : étude d'un modèle discret}
	
	\begin{enumerate}
	\item $u_{1}=4-\dfrac{4}{u_{0}}= 4 - \dfrac{4}{4}=4-1=3$.

Le nombre de perches-soleil au 1\up{er} janvier 2026 donné par ce modèle sera de \np{3000}.

\item 
	\begin{enumerate}
		\item La fonction $h$ est dérivable sur $] 0;~+\infty[$ et pour tout réel strictement positif $x$ on a :

$h'(x) = 0 -4 \times \dfrac{-1}{x^2}= \dfrac{4}{x^2} > 0$

Donc $h'$ est une fonction à valeurs strictement positives sur $]0~;~+\infty[$, cela implique que la fonction $h$ est strictement croissante sur l'intervalle $]0~;~+\infty[$.
		\item On remarque que la fonction $h$ introduite dans la question est la fonction de récurrence de la suite, c'est-à-dire que pour tout entier naturel $n$, on a : \quad $h(u_n) = u_{n+1}$.

Pour tout entier naturel $n$, on pose l'affirmation $P_n$ :\quad \og $2 \leqslant u_{n+1} \leqslant u_{n} \leqslant 4$.\fg{}

\textbf{Initialisation : } pour $n = 0$, on a, d'après la définition de la suite $u_0 = 4$ et d'après la question précédente $u_1 = 3$.

Donc \quad $2 \leqslant u_1 \leqslant u_0 \leqslant 4$. \quad Donc l'affirmation $P_0$ est vraie.

\textbf{Hérédité :} Soit $n\in \N, \; n \geqslant 0$ et $2 \leqslant u_{n+1} \leqslant u_{n} \leqslant 4$ vraie (c'est l'hypothèse de récurrence $P_n$)

Par croissance de la fonction $h$ démontrée à la question \textbf{2. a.} les images de ces quatre nombres sont rangées dans le même ordre que leurs antécédents, soit

$h(2) \leqslant h(u_{n+1}) \leqslant h(u_n) \leqslant h(4)$.

Or $h(2) = 4 - \dfrac42  = 4 - 2 = 2, \quad h(u_{n+1}) = u_{n+2}, \quad h(u_n) = u_{n+1}$ et \quad $h(4) = 4 - \dfrac 44 = 4 - 1 = 3$.

On a donc $2 \leqslant u_{n+1} \leqslant u_{n+2} \leqslant 3$ et comme $3 < 4$, on a finalement 

$2 \leqslant u_{n+1} \leqslant u_{n+2} \leqslant 4$ : $P_{n+1}$ est vraie.
















\textbf{Conclusion : } la relation est vraie au rang 0 et si elle est vraie au rang $n \geqslant 0$, elle l'est aussi au rang $n+1$ : d'après le principe de récurrence :
\[\forall n \in \N,\quad 2 \leqslant u_{n+1} \leqslant u_{n} \leqslant 4.\]

		\item D'après la question précédente, on a :

		\begin{itemize}
			\item $\forall n \in \N,\quad 2\leqslant u_n \leqslant 4$ : \quad la suite est minorée par 2 et majorée par 4 ;
			\item $\forall n \in \N,\quad u_{n+1} \leqslant u_{n}$:\quad la suite est décroissante;
		\end{itemize}

La suite est donc décroissante et minorée, elle est donc convergente, vers une limite $\ell~\geqslant~2$.



		\item La suite $u$ est définie par récurrence, sa fonction de récurrence $h$  est continue car dérivable sur $]0~;~+\infty[$, et la suite est  convergente. D'après le théorème du point fixe, on en déduit que la limite $\ell$  de la suite $(u_n)$ doit vérifier l'équation $h(\ell) = \ell \iff 4 - \dfrac{4}{\ell} = \ell$.











En multipliant chaque membre par $\ell \ne 0$, on obtient $\ell^2  = 4\ell  - 4 \iff \ell^2 - 4\ell + 4 = 0 \iff (\ell - 2)^2 = 0 \iff \ell - 2 = 0 \iff \ell = 2$.

L'unique solution de l'équation étant 2, on a donc $\displaystyle\lim_{n \to + \infty} u_n = 2$.




		\item Ce modèle prévoit une stabilisation de la population à \np{2000} perches-soleil, et ne prévoit donc pas une élimination à long terme de l'espèce envahissante.
	\end{enumerate}

	\item \begin{enumerate}
		\item Soit $s$ un réel appartenant à l'intervalle $]2~;~4[$.

Le script qui renvoie, après exécution, le plus petit entier $n$ tel que $u_{n}<s$ est :

		\begin{center}
			\begin{ttfamily}
			\begin{tblr}{colspec={X[6cm]},hline{1,8}={solid},
			vlines,rows={rowsep=0pt}}
			def population(s) :\\
			\quad u = 4\\
			\quad n = 0\\
			\quad while u >= s :\\
			\quad \quad u = 4 - 4/u\\
			\quad \quad n = n + 1\\
			\quad return n
			\end{tblr}
		\end{ttfamily}
		\end{center}

La variable \texttt{u} va stocker la valeur d'un terme de la suite, tandis que la variable \texttt{n} va stocker la valeur de l'indice correspondant.

Pour que la boucle \texttt{while} s'arrête dès que $u_n < s$, il faut qu'elle continue tant que cette condition n'est pas réalisée, c'est-à-dire tant que \texttt{u} contient une valeur qui n'est pas strictement inférieure à $s$.

Et tant que c'est le cas, les deux instructions à répéter sont : de remplacer la valeur stockée dans \texttt{u} par la valeur suivante, en appliquant la relation de récurrence, et donc d'augmenter de 1 la valeur de l'indice stocké dans \texttt{n}, pour garder la correspondance.
		\item L'algorithme donne :

$\bullet~$ $u_9 = 2,2$ et 

$\bullet~$ $u_{10} \approx 2,182 <  2,2$, donc la commande \texttt{population(2.2)} renvoie 10.

La population sera inférieure à \np{2200} perches-soleil à partir de 2035 (2025 + 10).
	\end{enumerate}
	\end{enumerate}

	\textbf{Partie B : étude d'un modèle continu}
	
	\begin{enumerate}
		\item $(E) \quad y'+y=2$ \quad soit \quad $y'=-y+2$.

		Cette équation différentielle est de la forme \quad $y'=ay+b$ \quad avec $a=-1$ et $b=2$.

		Les solutions de l'équation $(E)$ sont donc les fonctions définies sur $\R$ de la forme :

		$f(t) = k \e^{-t} - \dfrac{2}{-1}= k\e^{-t} +2$ \quad avec $k \in \R$.


		\item $p$ est solution de l'équation $(E)$ donc il existe un réel $k$ tel que $p(t) = k\e^{-t} +2$

		On sait que $p(0)=4$ donc : \quad $\begin{aligned}[t]
			p(0) = 4 &\iff 4= k \e^{0} +2 \\
		&\iff 2 = k
		\end{aligned}$

	L'expression de la fonction $p$ sur l'intervalle $[0;+\infty[$ est donc bien : \quad $p(t)=2 \e^{-t}+2$.
		\item Comme $\lim\limits_{t\to +\infty} \e^{-t} = 0$, par limite de produit par 2 puis limite de somme $\lim\limits_{t\to +\infty} 2 \e^{-t} + 2 = 2$

Donc ce modèle ne prévoit pas non plus une élimination à long terme de l'espèce envahissante, il prévoit une stabilisation à \np{2000} perches-soleil. Ce qui est cohérent avec le modèle étudié à la partie \textbf{A}.
	\end{enumerate}
